This study shows that modifying the training objective is an effective and
controlled way to influence model behavior in sparse lesion segmentation.
Instead of increasing architectural complexity, the proposed approach
changes how supervision is distributed across lesions, which makes the
effect of the loss design directly interpretable.

\paragraph{Why standard losses miss small lesions.}
The probability analysis shows that the baselines do not miss small lesions
by a narrow margin: for most of them, the network produces no response at
all. A lower decision threshold therefore cannot recover these lesions, and
the omission has to be addressed during training. The comparison with the
Tversky-based losses locates the cause. Tversky and FocalTversky penalize
false-negative voxels more heavily than DiceCE and recover part of the
lesions that DiceCE misses, but they remain behind CATMIL in the cases with
many small lesions. Because Tversky weighting applies to all false-negative
voxels alike, a lesion of a few voxels is still a negligible fraction of
the loss. Only when each lesion is treated as a unit, as in the CAT and MIL
terms, does every lesion contribute to the objective regardless of its
size. The gain of CATMIL therefore comes from lesion-level supervision
rather than from a stronger penalty on false negatives, and on MSLesSeg it
is obtained without loss of voxel-level overlap or boundary accuracy. The
ablation in \cref{appendix:catmil-components} is consistent with this
interpretation: the MIL term alone already raises small-lesion recall,
whereas the combination with the CAT term is needed to keep DSC and HD95.

\paragraph{The trade-off and its control.}
The greater sensitivity of CATMIL comes with more false-positive components
and a lower lesion-wise F1. These components are small, lie mostly close to
true lesions, and contribute little false-positive volume. They reflect the
same responsiveness to weak signals that allows CATMIL to detect small
lesions, extended to ambiguous patterns that resemble lesions. Because they
are small, a component-size filter removes most of them while the
sensitivity gain is retained, so the trade-off can be adjusted after
training without retraining the model. In addition, the MIL term requires
only one confident voxel per lesion and does not enforce full lesion
coverage, which explains why some recovered lesions are detected but only
partly delineated.

\paragraph{Clinical implications.}
In multiple sclerosis, the number of lesions and the appearance of new
lesions on follow-up scans are central to diagnosis and to the monitoring
of disease activity. A small lesion that is never detected cannot be
counted, whereas a false-positive component can be dismissed by a reader
when the segmentation is reviewed. For this setting, reducing missed
lesions at the cost of additional small false positives is a favorable
direction for the trade-off, particularly when the automatic segmentation
supports rather than replaces the radiologist. The component-size filter
then allows the operating point to be chosen for the intended use: a small
filter for screening and follow-up, where sensitivity matters most, and a
larger filter where a cleaner lesion map is needed.

\paragraph{Limitations.}
Several limitations should be noted. First, the evaluation is restricted to
multiple sclerosis, and the applicability of the proposed objective to
other brain pathologies remains to be validated; all experiments also use
the nnU-Net framework, so the effect of CATMIL on other architectures
remains to be tested. Second, the 3D-MR-MS dataset was evaluated with the
same training configuration as the main experiments to keep the comparison
consistent, so some hyperparameters may not be optimal for its
characteristics; the false-positive cost of CATMIL is larger on this
dataset, and a larger filter is required to control it. The filter size is
thus chosen after training and depends on the data, rather than being
controlled within the objective. Third, the connected components used by
the loss are computed within each training patch, so a lesion cut by the
patch border is weighted according to the part visible in the patch rather
than its full size. Fourth, the test sets are small (six patients each),
so the per-case observations describe a consistent within-dataset pattern
rather than a population-level estimate. Finally, lesions of only a few
voxels and lesions with low contrast or diffuse boundaries remain difficult
for every loss: the objective increases the influence of small lesions
during training, but it cannot compensate for missing evidence in the
image.

\paragraph{Future work.}
These limitations suggest several directions. The objective should be
evaluated on other conditions with small focal lesions, such as brain
metastases, satellite glioma lesions, or microbleeds, where missed lesions
are also clinically important, and with other segmentation architectures.
Methodologically, constraints on lesion coverage, boundary consistency, or
structural connectivity could reduce fragmented predictions and isolated
false positives, and could move the control of the precision trade-off
from post-processing into the objective itself. A natural candidate is a
penalty on predicted components that overlap no reference lesion, as in
BiCC~\cite{bouteille_bicc_2026}: it acts on exactly the error that CATMIL
introduces, whereas the MIL term acts on the error that CATMIL removes, so
the two are complementary. Finally, an analysis of the
learned features could show how the objective changes the internal
representation of small lesions, beyond the output-level effects reported
here.
